\bibitem{support-spectrum-Caraiani2012}
A. Caraiani, \emph{Local-global compatibility and the action of
monodromy on nearby cycles}, Duke Mathematical Journal
\textbf{161} (2012), 2311--2413.  Theorems~1.1--1.2,
p.~2312: full Weil--Deligne compatibility away from the
coefficient prime and temperedness for regular algebraic
conjugate self-dual cuspidal representations over CM fields.
\url{https://doi.org/10.1215/00127094-1723706}.
Author's copy: \url{https://www.ma.imperial.ac.uk/~acaraian/papers/lgc1.pdf}.

\bibitem{support-spectrum-Caraiani2014}
A. Caraiani, \emph{Monodromy and local-global compatibility for
$l=p$}, Algebra \& Number Theory \textbf{8} (2014),
1597--1646.  Theorem~1.1, pp.~1597--1598, including the
monodromy operator and without a weight bound relative to $l$.
\url{https://doi.org/10.2140/ant.2014.8.1597}.
\url{https://msp.org/ant/2014/8-7/ant-v8-n7-p02-s.pdf}.

\bibitem{support-spectrum-Saito}
T. Saito, \emph{Modular forms and $p$-adic Hodge theory},
Inventiones Mathematicae \textbf{129} (1997), 607--620;
\emph{Weight-monodromy conjecture for $\ell$-adic representations
associated to modular forms. A supplement}, in
B. B. Gordon et al., eds., \emph{The Arithmetic and Geometry of
Algebraic Cycles}, NATO Science Series C \textbf{548} (2000),
427--431.  The supplement's theorem on p.~427 and corollary
on p.~428 give the full Frobenius-semisimplified Weil--Deligne
comparison at the coefficient prime, for newforms of weight
at least two.  The supplement is part of the cited result.
\url{https://doi.org/10.1007/s002220050175}.
Author's copy of the supplement:
\url{https://www.ms.u-tokyo.ac.jp/~t-saito/pp/weightmonodromy.pdf}.
The hypotheses and the role of this supplement are also
recorded in the author's \emph{Hilbert modular forms and
$p$-adic Hodge theory}, Compositio Mathematica
\textbf{145} (2009), 1081--1113, Introduction and
Theorems~2.1--2.2; \url{https://doi.org/10.1112/S0010437X09004175}.

\bibitem{support-spectrum-newvector}
H. Jacquet, I. I. Piatetski-Shapiro and J. Shalika,
\emph{Conducteur des repr\'esentations du groupe lin\'eaire},
Mathematische Annalen \textbf{256} (1981), 199--214,
Theorem~5.1, p.~211.  A generic local representation of
conductor exponent $a$ has a nonzero $K_1(\mathfrak p^a)$-fixed
vector.  The proof uses H. Jacquet's correction to the
essential-vector argument, \emph{A correction to Conducteur
des repr\'esentations du groupe lin\'eaire}, Theorem~1.
\url{https://www.math.columbia.edu/~hj/conducteur.pdf}.
\url{https://www.math.columbia.edu/~hj/CorrectionToConducteur.pdf}.

\bibitem{support-spectrum-BorelSerre}
A. Borel and J.-P. Serre, \emph{Corners and arithmetic groups},
Commentarii Mathematici Helvetici \textbf{48} (1973),
436--491, Theorem~9.3, p.~475.  At a fixed neat level the
Borel--Serre compactification is a compact manifold with
corners, yielding a finite cell complex; the complex lifts
to every finite covering.
\url{https://doi.org/10.1007/BF02566134}.
\url{https://www.e-periodica.ch/cntmng?pid=com-001%3A1973%3A48%3A%3A32}.

\bibitem{support-spectrum-Clozel}
L. Clozel, \emph{Motifs et formes automorphes: applications
du principe de fonctorialit\'e}, in L. Clozel and J. S. Milne,
eds., \emph{Automorphic Forms, Shimura Varieties, and
$L$-Functions}, Vol.~I, Academic Press (1990), 77--159.
Theorem~3.13, p.~112; Lemma~3.14, p.~114; and
Lemma~3.15, p.~121.  These supply rationality and scalar
conjugation, nonzero cohomology with the indicated algebraic
coefficients, and the injection of cuspidal into ordinary
cohomology.  The editor's digitized primary volume is at
\url{https://www.jmilne.org/math/Books/AA1988a.pdf}.

\bibitem{support-spectrum-BLGGT}
T. Barnet-Lamb, T. Gee, D. Geraghty and R. Taylor,
\emph{Potential automorphy and change of weight},
Annals of Mathematics \textbf{179} (2014), 501--609,
Theorem~2.1.1.  Parts~(2)--(4) give the compatible
representations, labeled Hodge--Tate weights and the
crystalline conclusion at an unramified local component.
Its Iwahori restriction in part~(4) is sufficient for that
last application; arbitrary ramified coefficient-prime
components are handled by the separately cited full
compatibility theorems.
\url{https://doi.org/10.4007/annals.2014.179.2.3}.
\url{https://www.ma.imperial.ac.uk/~tsg/Index_files/potentialautomorphy.pdf}.

\bibitem{support-spectrum-Deligne}
P. Deligne, \emph{Formes modulaires et repr\'esentations
$\ell$-adiques}, S\'eminaire Bourbaki, Expos\'e~355,
Lecture Notes in Mathematics \textbf{179} (1971), 139--172;
\emph{La conjecture de Weil: I}, Publications
Math\'ematiques de l'IH\'ES \textbf{43} (1974), 273--307.
The modular-form construction, Hecke integrality, and
the Weil theorem give the good-prime roots in (SP14),
with absolute value $p^{(k-1)/2}$ at every complex embedding.
\url{https://www.numdam.org/item/SB_1968-1969__11__139_0.pdf}.
\url{https://www.numdam.org/articles/10.1007/BF02684373/}.

\bibitem{support-spectrum-source}
L. Alp\"oge, pinned source mirror \path{jsmorph/rh}, revision
\newline\texttt{ecd09ff9b4aaf6c08f2b5ef17dc49ca9648a77de}.
Under \path{Development/AutomorphicArithmetic/AutomorphicVarieties/}:
\path{AutomorphicForms/BlochKatoSelmerVanishes.lean},
66--114; \path{AutomorphicForms/CongruenceNumberRankThree.lean},
56--92 and 107--140; \path{AutomorphicForms/Newform.lean},
21--47; \path{AutomorphicForms/TemperedLocalComponents/CountK.lean},
41--77; \path{AutomorphicForms/TemperedLocalComponents/CohBoundFree.lean},
69--99; and
\path{AutomorphicForms/TemperedLocalComponents/CongruenceNumberRankThree.lean},
14--63.  The certificate authenticates these source interfaces
and their original quantifiers; it is not a Lean rebuild.

\bibitem{support-spectrum-orbit-source}
The same pinned source revision:
\path{HeightAndVanishing/Counts.lean}, 9--31 and 53--58;
\path{HeightAndVanishing/GdFacts.lean}, 22--26 and 65--101;
\path{UnitaryThreeGroups/KatoUse.lean}, 36--64;
\path{UnitaryThreeGroups/OrbEquiv.lean}, 71--109 and 111--173;
\path{UnitaryThreeGroups/CorGSumLog.lean}, 6--73; and
\path{UnitaryThreeGroups/CountI.lean}, 22--136.
The fixed field in the last application is chosen from the
fixed primitive character before the host parameters vary.

\bibitem{support-spectrum-case-one-source}
The same pinned source revision:
\path{HeightAndVanishing/DirichletPrimeToLPartMulLPowerPart.lean},
148--192 and 206--219;
\path{HeightAndVanishing/BchiFacts.lean}, 7--10 and 22--41;
\path{HeightAndVanishing/GdOfHost.lean}, 37--94; and
\path{UnitaryThreeGroups/CountI.lean}, 22--136.
These give the original arithmetic family, fixed exceptional
set, ramification verification and kernel-orbit summation.
